\newpage
\section{CODESYS codificación manual}
Los Niveles 0 y 1 se han diseñado en Stateflow y el Nivel 2 en Simulink como diagrama de bloques discretos, sin embargo, 
para la implementación en hardware real, se debe traducir la lógica desarrollada a un programa en CODESYS utilizando una mezcla de los lenguajes
diagrama de funciones secuenciales (SFC) y texto estructurado (ST). Este proceso implica mapear 
cada bloque funcional a funciones y estructuras de control equivalentes en SFC y ST, asegurando que la lógica de control se mantenga 
intacta durante la transición al entorno de ejecución real. Puede verse a continuación la estructura resultante:
\begin{figure}[H]
    \centering
    \includegraphics[width=0.35\textwidth]{codesys/arbol_codesys.png}
    \caption{Árbol de componentes del programa en CODESYS}
\end{figure}

Los dos POUs principales corresponden a los Niveles 0 y 1, mientras que el Nivel 2 se implementa como bloques de función 
desarrollados en Matlab con comportamiento equivalente al lenguaje ST y ejecutados
en tiempo discreto. Estos dos POUs se ejecutan de forma cíclica e independiente. En un sistema con múltiples CPUs podrían ejecutarse en paralelo; sin embargo, en CODESYS, debido a la utilización de un PLC virtual con una única CPU, ambas tareas se ejecutan de forma secuencial, alternando su ejecución en cada ciclo de planificación. 
A pesar de la limitación, se ejecuta a tal velocidad que, a fines prácticos, el comportamiento es equivalente al de una ejecución paralela.

\subsection{Control regulatorio en simil ST}
Se rediseña el modelo del Nivel 2 usando la misma lógica pero representándola con bloques de función de Matlab (simil ST):
\begin{figure}[H]
    \centering
    \includegraphics[width=0.55\textwidth]{codesys/control_carro_funciones2.png}
    \caption{Diagrama de control del carro con funciones de Matlab}
\end{figure}
\begin{figure}[H]
    \centering
    \includegraphics[width=0.65\textwidth]{codesys/control_izaje_funciones2.png}
    \caption{Diagrama de control de izaje con funciones de Matlab}
\end{figure}
Se agregan varios bloques \textit{Delay} para solucionar problemas de loops algebráicos en la simulación.
El código de cada controlador queda de la siguiente forma:
\begin{figure}[H]
    \centering
    \begin{subfigure}{0.49\textwidth}
        \centering
        \includegraphics[width=\linewidth]{codesys/matlab_function_nivel2_carro.png}
        \caption{Control del carro simil ST}
    \end{subfigure}
    \hfill
    \begin{subfigure}{0.49\textwidth}
        \centering
        \includegraphics[width=\linewidth]{codesys/matlab_function_nivel2_izaje.png}
        \caption{Control de izaje simil ST}
    \end{subfigure}
    \caption{Código de los controladores del Nivel 2 en simil ST}
\end{figure}

\subsection{Control supervisor en CODESYS}
En el control supervisor se pasan los bloques de Stateflow a SFC, manteniendo el mismo número de estados y lógica de transición. 
De este proceso se obtiene un conjunto de varios bloques en paralelo:
\begin{figure}[H]
    \centering
    \includegraphics[width=1\textwidth]{codesys/codesys_supervisor11.png}
    \caption{Diagrama general de control supervisor en CODESYS}
\end{figure}

El autómata principal se desarrolla de forma similar, y se encarga de alternar entre los modos de reposo, 
manual, automático, errores y recuperación. 
\begin{figure}[H]
    \centering
    \begin{subfigure}{1\textwidth}
        \centering
        \includegraphics[width=\linewidth]{codesys/codesys_supervisor12.png}
    \end{subfigure}
    \begin{subfigure}{1\textwidth}
        \centering
        \includegraphics[width=\linewidth]{codesys/codesys_supervisor13.png}
    \end{subfigure}
    \caption{Diagrama del autómata principal en CODESYS}
\end{figure}

Para el modo automático se realizan algunos cambios. En Stateflow, al entrar al modo automático se pasa por una \textit{junction} para decidir si la ejecución comienza desde cero o si se retoma desde el paso interrumpido. En SFC esto no es posible, ya que el diagrama debe iniciar obligatoriamente en un estado, por lo que se utiliza un estado vacío.
Además, al finalizar el movimiento de \textit{homing}, en Stateflow la transición de salida del modo automático conduce directamente al modo reposo. En cambio, en SFC es necesario incluir una transición o un salto luego del estado para que el programa compile correctamente. Por este motivo, se agrega una transición con condición \textit{FALSE}, la cual nunca se cumple.
En la siguiente imágen se puede ver la implementación utilizada:
\begin{figure}[H]
    \centering
    \includegraphics[width=0.45\textwidth]{codesys/codesys_supervisor15.png}
    \caption{Diagrama del modo automático en CODESYS}
\end{figure}


Sin modificaciones notorias, se replica el bloque encargado de la lógica de los Twistlocks. Como es un SFC cíclico, el último paso regresa al inicial al cumplir su transición:
\begin{figure}[H]
    \centering
    \includegraphics[width=0.75\textwidth]{codesys/codesys_supervisor4.png}
    \caption{Diagrama de la lógica de los Twistlocks en CODESYS}
\end{figure}

De la misma forma, se utilizan las entradas de los sensores para calcular la posición discreta del carro y actualizar el perfil de 
obstáculos, información que es fundamental para la generación de trayectorias en modo automático. También maneja la lógica del
escaneo del muelle:
\begin{figure}[H]
    \centering
    \includegraphics[width=0.6\textwidth]{codesys/codesys_supervisor14.png}
    \caption{Diagrama del bloque indices perfil en CODESYS}
\end{figure}

Los bloques de \textit{Suavizado velocidad}, \textit{Velocidad límites} y \textit{Watchdog heartbeat} se mantienen prácticamente iguales, desarrollados en ST directamente debido a su sencillez.

\subsection{Control de protección en CODESYS}
Para el nivel de protección se mantiene el mismo diseño, con dos procesos paralelos, uno para la máquina de estados de seguridad y 
otro para el watchdog. La lógica de ambos procesos es prácticamente la misma que en Stateflow, con pequeñas modificaciones para adaptarse a la 
sintaxis y estructura de SFC.
\begin{figure}[H]
    \centering
    \includegraphics[width=0.5\textwidth]{codesys/codesys_seguridad1.png}
    \caption{Diagrama general del control de protección en CODESYS}
\end{figure}

\begin{figure}[H]
    \centering
    \includegraphics[width=1\textwidth]{codesys/codesys_seguridad22.png}
    \caption{Diagrama de la máquina de estados de seguridad en CODESYS}
\end{figure}

\begin{figure}[H]
    \centering
    \includegraphics[width=0.83\textwidth]{codesys/codesys_seguridad3.png}
    \caption{Diagrama del watchdog en CODESYS}
\end{figure}

\subsection{Comunicación OPC Simulink/CODESYS}
Para poder comunicar el PLC virtual de CODESYS con Simulink, se usa el protocolo de comunicación industrial \texttt{OPC-UA}. La versión de CODESYS utilizada (CODESYS V3.5 SP21 Patch 2) cuenta con la posibilidad de montar un servidor OPC-UA y de seleccionar qué variables del programa se quieren exponer a través de este protocolo. Por otro lado, Matlab R2025b cuenta con un cliente OPC-UA que permite conectarse a este servidor y leer/escribir las variables expuestas en tiempo real.

El único ajuste de seguridad que se modifica es el de \textit{User allow anonymous},
lo que permite a Simulink conectarse al servidor OPC-UA de CODESYS sin problemas, y leer/escribir las variables expuestas en tiempo real. Sin embargo, lo hace de forma anónima y sin seguridad alguna, lo cual no es recomendable para un entorno real, pero para fines de simulación y pruebas es suficiente.

Debido a la naturaleza de la comunicación OPC, cada lectura o escritura de una variable implica un proceso de comunicación 
cliente-servidor que introduce un retardo significativo. Este retardo se vuelve especialmente problemático cuando se intenta 
realizar lecturas/escrituras a alta frecuencia, lo cual es necesario para el control supervisor. OPC es un estandar de tipo TCP/IP, el cuál está diseñado para sistemas SCADA, con retardos de comunicación de $500$ ms o más, y no para sistemas de control en tiempo real. 

Para solucionar este problema, se implementan bloques de comunicación que se ejecutan a diferentes tiempos de muestreo. 
El control supervisor y de protección se ejecutan a $500$ ms y para no tener errores 
con el heartbeat del watchdog, su variable corre al tiempo de muestro original de $20$ ms. De esta forma, se reduce la cantidad
de lecturas/escrituras por segundo, disminuyendo la carga en la red y permitiendo Simulink pueda correr a $1$s por segundo de simulación.

Con el uso de los bloques de comunicación con mayor tiempo de muestreo, se mejora significativamente la fluidez de la simulación, manteniendo una comunicación efectiva entre Simulink y CODESYS. Además, se comprueba de forma satisfactoria que el controlador responde de la misma manera que el modelo en Simulink.

A continuación, los bloques de escritura/lectura de las variables de control supervisor y protección y una captura de la simulación 
con la comunicación OPC-UA funcionando, respectivamente:

\begin{figure}[H]
    \centering
    \includegraphics[width=0.79\textwidth]{codesys/opc_ua_1.png}
    \caption{Diagrama de bloques de comunicación OPC-UA para el control supervisor en Simulink}
\end{figure}
\begin{figure}[H]
    \centering
    \includegraphics[width=0.79\textwidth]{codesys/opc_ua_0.png}
    \caption{Diagrama de bloques de comunicación OPC-UA para el control de protección en Simulink}
\end{figure}
\begin{figure}[H]
    \centering
    \includegraphics[width=0.8\textwidth]{codesys/codesys_interfaz.png}
    \caption{Simulación con comunicación OPC-UA en funcionamiento.}
\end{figure}


\section{Generación automática de código ST}
Simulink ofrece una herramienta llamada PLC Coder, que tiene la capacidad de traducir automáticamente diagramas Stateflow a ST bajo 
la norma IEC 61131-3. Analizando el código ST resultante de esta autogeneración, se puede observar que el traductor 
implementa el autómata de forma ``plana'' y monolítica. Para lograrlo, recurre a una estructura de control \texttt{CASE} global que evalúa una variable entera 
representativa del estado actual, combinada con una densa red de sentencias \texttt{IF/ELSIF} anidadas para resolver las transiciones y acciones correspondientes.
A continuación se muestran fragmentos del código autogenerado para el control de seguridad, evidenciando 
la complejidad y falta de claridad que presenta:

\begin{figure}[H]
    \centering
    \includegraphics[width=0.85\textwidth]{codesys/st_generado_2.png}
    \caption{Fragmento de código ST autogenerado por Simulink PLC Coder en CODESYS}
\end{figure}


Además de la complejidad visual, el código autogenerado es un archivo .ST para cada nivel, por lo que para su uso en CODESYS es necesario 
dividir manualmente cada archivo en varios bloques de función. También se debe crear y llamar una instancia del bloque de función de cada nivel en el programa principal,
ya que este es el que se comunica con Simulink a través de OPC-UA, conteniendo las variables de entrada/salida que se exponen a través del protocolo. El programa
principal y la estructura de bloques de función resultante del código autogenerado se muestra a continuación:

\begin{figure}[H]
    \centering
    \begin{subfigure}{0.35\textwidth}
        \centering
        \includegraphics[width=\linewidth]{codesys/arbol_codesys_st.png}
        \caption{Estructura de bloques de función}
    \end{subfigure}
    \hspace{0.05\textwidth}
    \begin{subfigure}{0.35\textwidth}
        \centering
        \includegraphics[width=\linewidth]{codesys/codesys_main_st.png}
        \caption{Programa principal}
    \end{subfigure}
    \caption{Estructura de bloques de función y programa principal resultante del código autogenerado}
\end{figure}


En teoría a nivel de eficiencia computacional y carga de CPU, el rendimiento del código autogenerado es comparable al de la implementación manual. 
Los compiladores modernos de los PLC optimizan ambas estructuras a un código máquina de ejecución sumamente rápida, por lo que el 
tiempo de ciclo del autómata no se ve significativamente alterado por la elección de uno u otro método.

Sin embargo, la divergencia crítica entre ambos enfoques radica en la legibilidad, la mantenibilidad y el diagnóstico de fallas en el 
entorno industrial. El código ST autogenerado carece de trazabilidad visual; resulta ser una estructura de difícil interpretación para 
los operarios de planta. 

\begin{figure}[H]
    \centering
    \includegraphics[width=0.7\textwidth]{codesys/codesys_interfaz_st.png}
    \caption{Simulación con comunicación OPC-UA en funcionamiento usando código ST autogenerado}
\end{figure}